\subsection{Detection by the actual endpoint period}

\begin{proposition}[DP55: endpoint detection and a nonzero principal-level period]\label{DP55}
Assume the weak conjugation hypotheses and specified endpoint weights and
level of DP52(c). Let $S,S'$ be the actual point and row data, and
$\mathcal E=\operatorname{ptEv}$ their point evaluation object. Put
\[
 \mathcal F=\operatorname{span}_{\mathbf C}
    \{\mathcal E.\operatorname{emb}(\phi):
                    \phi\in S.\operatorname{EqFun}\}.
\]
This is precisely the source's $\mathcal E.\operatorname{Fi}$.
\begin{enumerate}[label=(\alph*),leftmargin=*]
\item The source row representing map
$J:\mathbf C^{\dim_K S'.\operatorname{EqFun}}\to\mathcal F$
is conjugate-linear and surjective under these weaker hypotheses.
The actual weakened endpoint period satisfies
\[
 \operatorname{per}(F,z)=\langle Jz,F\rangle_{\mathrm{cl}},
\]
where the classwise form is positive definite on $\mathcal F$.
Thus every nonzero $F\in\mathcal F$ has a row vector $z$ with
\[
 Jz=F,\qquad
 \operatorname{per}(F,z)=\langle F,F\rangle_{\mathrm{cl}}>0.
\]
This requires neither a free class action nor a change of the given level.
\item Add an auxiliary place $P$ above an odd rational prime $\ell$
unramified in the CM field $E$, outside the old finite endpoint support.
Set its level subgroup equal to $\{1\}$, with trivial character, retaining
the other level data.
The decorated-lattice action is then free. For every dominant endpoint
coefficient weight, the point and row spaces are nonzero and the actual
endpoint period has a pair with value $1$.
This applies to the actual type $\{1,2,3\}$ of $\mathbf Q(\zeta_7)$
at every admissible rebuilt even lift count.
\item A vector belongs to $\mathcal F$ exactly when its value at every
class lies in the coefficient-evaluation span and is invariant under the
class stabilizer. The source's adelic identification additionally requires
the explicit inputs $\operatorname{DecoPtSurj}$ and
$\operatorname{IsLevOf}$; that theorem has no
$\operatorname{Fields.Ok}$ argument.
\end{enumerate}
The proposition supplies no quantitative bound on the chosen row
coordinates or the cost of the auxiliary level, and does not construct
a nonzero image of the prescribed iterated theta kernel.
\end{proposition}

\begin{proof}
The source paths in this proof begin with
\path{Formalism/LocalToGlobal/Columns/}.
The actual point and row representations are respectively the round
and square coefficient representations, with inverse finite-level
characters. They have the same decorated-lattice set and identified
acting groups.
Their perfect bilinear coefficient pairing is
$B=\operatorname{pairBK}_\lambda$.
Its exact definition is the restriction of
\[
 B_\lambda(f,g)=
 \frac{\sum_\alpha f_\alpha g_\alpha\,\alpha!}
 {\prod_\tau\sum_\alpha
       (\operatorname{extremal}_{\lambda_\tau})_\alpha^2\,\alpha!}.
\]
The actual theorems \path{pairingBInvariant_holds} and
\path{pairingBNondegenerate_holds} in \path{CoefModB.lean},
188--203, give
\[
 B(\rho_{\mathrm{pt}}(g)x,\rho_{\mathrm{row}}(g)y)=B(x,y)
\]
and nondegeneracy. The inverse finite characters cancel.
Thus \path{TiedPer.lean}, \path{perDua}, 232--239, constructs the
required coefficient duality with no common-quadratic-subfield hypothesis.
Here \path{PairBNondeg} at \path{CoefModB.lean}, 88--90, is only the
proposition; the cited theorem at 193--203 is its proof.

Let $H$ be the point coefficient hermitian form.
For $\phi\in S.\operatorname{EqFun}$, perfectness gives a unique
row value $\psi(d)$ at each decorated lattice such that
\[
 B(x,\psi(d))=H(\phi(d),x)\quad\text{for all }x.
\]
This assignment is row-equivariant, since
\[
 \begin{aligned}
 B(x,\psi(gd))
 &=H(\rho_{\mathrm{pt}}(g)\phi(d),x)\\
 &=H(\phi(d),\rho_{\mathrm{pt}}(g)^{-1}x)\\
 &=B(\rho_{\mathrm{pt}}(g)^{-1}x,\psi(d))\\
 &=B(x,\rho_{\mathrm{row}}(g)\psi(d)).
 \end{aligned}
\]
Uniqueness in the second variable proves the assertion.
This is the complete argument of \path{GSD.Dua.partner} at
\path{TiedPer.lean}, 104--117; it uses DP52's explicit point-goodness
invariance.

Choose a $K$-basis $(b_a)$ of row equivariant functions and write
$\psi=\sum_a c_ab_a$.
For chosen point-class representatives $d_c$, the representing map is
\[
 (Jz)(c)=\sum_a\overline{z_a}\,
  \mathcal E.\operatorname{sharp}
  \bigl(x\longmapsto\sigma_0(B(x,b_a(d_c)))\bigr).
\]
The coefficient Gram matrix is positive definite, hence invertible:
this is exactly the matrix identified in the proof
\path{RecordInst.ptEv_gram_unit}, 47--58.
The operation \path{sharp} represents the displayed functional in the
first, conjugate-linear slot of the coefficient inner product.
Putting $z_a=\sigma_0(c_a)$ gives
$Jz=\mathcal E.\operatorname{emb}(\phi)$.
Indeed, the pairings of the two sides with every evaluated coefficient
vector agree, both sides lie in the coefficient-evaluation span, and
the Gram form is nonsingular there.

The displayed formula makes $J$ additive and conjugate-linear.
If $F=\sum_i t_i\mathcal E.\operatorname{emb}(\phi_i)$ and
$Jz_i=\mathcal E.\operatorname{emb}(\phi_i)$, then
$J(\sum_i\overline{t_i}z_i)=F$.
Conversely every $Jz$ is coefficient-valued and stabilizer invariant by
row equivariance and pairing invariance; the averaging argument below
puts it in $\mathcal F$.
This proves exactly the stated surjectivity.
The corresponding full source proofs are
\path{jOf_mem_Fi}, 133--148,
\path{jOf_partner}, 150--162, and \path{jOf_onto}, 172--189,
in \path{TiedPer.lean}.
Its \path{per_jMC_onto}, 250--252, uses
$h_X:X.\operatorname{Ok}$ only through \path{ptOk}.
Replacing that argument by DP52's proved point-goodness certificate
preserves this entire construction.

The classwise form is
\[
 \langle F_1,F_2\rangle_{\mathrm{cl}}
 =\sum_c\frac{\langle F_1(c),F_2(c)\rangle_{\mathrm{coef}}}
                {|\operatorname{Stab}(d_c)|}.
\]
Every denominator is a positive finite integer.
Each coefficient form is positive definite on its evaluation span,
so the sum is positive for $F_1=F_2=F\ne0$.
The finite-sum period identity is
\path{Record.periodC_eq_ipCl}, 1580--1590, or
\path{RecordInst.periodC_eq_ipCl'}, 60--64, with DP52's weakened
goodness input.
Taking $Jz=F$ proves (a).

For the membership assertion, every evaluated equivariant function is
coefficient-valued and stabilizer invariant, and so is its complex span.
Conversely, express each invariant value $F(c)$ in the evaluated
$K$-coefficient basis, retaining its complex scalar coefficients.
Average each $K$-basis vector over the finite stabilizer of $d_c$,
dividing by the group order.
Their evaluated complex combination remains $F(c)$ by invariance.
Each averaged vector is a $K$-coefficient vector fixed by the stabilizer,
so it extends equivariantly to the orbit and by zero to other orbits.
Summing these finitely many $K$-valued orbit functions with the original
complex coefficients gives $F\in\mathcal F$.
This is the full argument in \path{TiedValFi.lean}, 16--133.

The same file's \path{fi_iff_FiC}, 239--255, identifies this space with
adelic automorphic coefficient values under an actual CM datum,
the goodness certificate, \path{DecoPtSurj} and \path{IsLevOf}.
The last condition, at 151--154, identifies the level subgroup with
\path{levU} and its character with \path{levTyp}.
The theorem constructs an adelic function from an equivariant function
in one direction; in the other, the level transformation law gives
stabilizer invariance and the preceding averaging proof applies.
There is no \path{Fields.Ok} argument in this general theorem.
This proves (c), with precisely its stated remaining inputs.

It remains to prove the principal-level assertion.
First, if $R$ is the ring of integers of a finite unramified extension
of $\mathbf Q_\ell$, with $\ell$ odd, the group
$1+\ell M_n(R)$ has no nonidentity torsion.
Otherwise a power would have prime order $q$ and could be written
$A=I+\ell^aC$, where $a\ge1$ and some entry of $C$ is a unit.
For $q\ne\ell$,
\[
 A^q\equiv I+q\ell^aC\pmod{\ell^{a+1}M_n(R)}\ne I.
\]
For $q=\ell$,
\[
 A^\ell\equiv I+\ell^{a+1}C
                  \pmod{\ell^{a+2}M_n(R)}\ne I.
\]
Indeed, intermediate binomial terms have valuation at least
$1+2a\ge a+2$, while the last has $a\ell\ge a+2$.
Both conclusions contradict prime order.

Now choose the auxiliary place in (b); only finitely many rational primes
are excluded by ramification or the old support.
If $(g,h)$ stabilizes a decorated lattice $d$ with lattice $N$, then
$gN=N$.
The standard hermitian form is positive at every complex embedding.
Its rational trace norm is positive definite on the underlying
$\mathbf Q$-space.
An isometry sends each member of a $\mathbf Z$-basis of $N$ into the
finite set of lattice points with its prescribed norm.
Thus the lattice-isometry group is finite and $g$ has finite order.

The exact decoration action in \path{StdLatUnimodular.lean}, 108--111,
is $\operatorname{be}_P(x)\mapsto
h_P\operatorname{be}_P(g^{-1}x)$.
At the new place $h_P=1$.
Stabilization and the \path{be_add}, \path{be_ker} fields of
\path{Deco.Ok}, 81--94, give
\[
 (g^{-1}-1)N\subseteq PN.
\]
After completion at $P$, the lattice is free over the unramified local
ring, and $P$ is generated by $\ell$.
The matrix of $g^{-1}$ therefore belongs to the torsion-free group just
considered. Hence $g=1$ locally, and so globally.
At every other level place, stabilization now gives
$h_w\operatorname{be}_w(x)=\operatorname{be}_w(x)$.
Surjectivity of \path{be_w}, another \path{Deco.Ok} field, forces
$h_w=1$. The full stabilizer is trivial.
The row and point actions have the same underlying action, so both
are free.

The class set is nonempty, by the actual standard-lattice construction
\path{DecoSet.base} at \path{StdLatUnimodular.lean}, line 650.
The endpoint rank is positive.
The coefficient module is nonzero: its tensor of highest-weight
extremal polynomials is a product of nonzero polynomials in disjoint
variable blocks. The full source membership and nonvanishing proofs
are \path{CoefColLemmas.lean}, 320--342.
For a free action, arbitrary values at class representatives extend
uniquely by $\phi(gd_c)=\rho(g)\phi(d_c)$;
this is \path{GSetLemmas.resFree}, 221--225.

Choose a class $c_0$, a nonzero coefficient $x$, and by perfectness
$y$ with $B(x,y)=1$.
Assign $x$ to the point value at $d_{c_0}$ and zero at other classes.
Assign $y$ to the row value at that same representative and zero elsewhere.
Expand the resulting row function in its chosen basis and apply
$\sigma_0$ to its coordinates, obtaining $z$.
The exact identity \path{Record.perCl_emb}, 331--337, has just one
contribution, with stabilizer order one:
\[
 \operatorname{per}(\mathcal E.\operatorname{emb}\phi,z)
                    =\sigma_0(B(x,y))=1.
\]
This proves (b) for the actual period.
DP52 verifies all weak assumptions for the genus-three type;
any odd auxiliary prime other than $7$, outside its finite support,
is unramified and may be used.
\end{proof}

The endpoint period now detects every nonzero vector in the exact source
coefficient space. An auxiliary principal level also gives a nonzero
vector and an actual period value $1$ for every dominant coefficient weight.
What remains is to identify a nonzero image of the intended iterated theta
family in this space, with its stated adelic, level and kernel ties, and to
control the row and level costs in the same global comparison.
A finite global theta-jet realization on a Jacobian does not alone supply
that identification. The proposition establishes neither a shorter
global lift chain nor minimality of the old chain.
